\documentclass{article}
\usepackage{amsmath}

\title{lec58}
\author{Dylan Tang}
\date{May 2026}

\begin{document}

\maketitle

\section{Review of 2-Sample Hypothesis Test, Efficiency, Robustness}
\begin{itemize}
    \item Recall that given $\bar{X}_C$ and $\bar{X}_T$, the two-sample test for no difference in means is given by
\[
    z = \frac{(\bar{X}_T - \bar{X}_C) - 0}{\frac{\sigma_T}{\sqrt{n_T}} + \frac{\sigma_C}{\sqrt{n_C}}}
\]
Then, we compare
\[
    2 \cdot P(Z > z) < \alpha
\]
\item \textbf{Efficiency} trusts the data, and tries to use as much data that the original setup provides as possible to reach a conclusion.
\item \textbf{Robustness} protects against bad assumptions, noise, or inaccurate data.
\end{itemize}

\section{Efficiency and Robustness with Regularizers}
\subsection{Ridge Regression}
Consider the objective of ridge regression:
\[
    \arg\min_w \sum_{i=1}^n|y-w^Tx|^2_2 + \lambda |w|^2_2
\]
This is equivalent to putting a \textbf{Gaussian prior} on the weights, and assuming 
\[
    w \sim N(0, \frac{1}{\lambda}I)
\]
which results in the optimal
\[
    \hat{w} = (X^TX + \lambda I)^{-1}X^Ty
\]
Ridge regression penalizes large $w_i$s with higher loss.
\begin{itemize}
    \item Therefore, a larger $\lambda$ means that we trust the prior more, and thus shrink $w_i$s to 0 $\implies$ robustness
    \item Conversely, a smaller $\lambda$ means that we trust the prior less, and thus allow for greater $w_i$s $\implies$ efficiency
\end{itemize}
\subsection{LASSO Regression}
\begin{itemize}
    \item LASSO regression also introduces penalties for large weights.
    \item However, rather than shrinking weights, LASSO regression pushes coefficients exactly to 0, and thus keeps contributions only from important weights.
\end{itemize}
The objective of LASSO regression is the following:
\[
    \arg\min_w \sum_{i=1}^n|y-w^Tx|^2_2 + \lambda |w|_1
\]
where
\[
    |w|_1 = \sum_{i=1}^k w_i
\]
The same concepts of efficiency and regularization also pertain here. A larger lambda means we trust LASSO prior more, and shrink more coefficients to 0. A smaller lambda means that we trust the LASSO prior less, and keep more coefficients.

\section{Distributional Robustness}
\begin{itemize}
    \item Recall that a regression model is given by 
    \[
        \hat{y} = \hat{m}X + \epsilon
    \]
    \item However, what if there is a small subset of data that are adversarial to our model, i.e. they lead our data to be inaccurate.
    \item Solution 1: Acquire more data to increase robustness
    \item Solution 2: Reduce the influence of one corrupted data point by converting to percentiles or ranks
        \begin{itemize}
            \item Reduces efficiency but increases robustness
        \end{itemize}
\end{itemize}

\newpage
\section{Three Types of  Priors}
\begin{itemize}
    \item \textbf{Smoothing prior} - Assume that $w_i$s are close to 0 - use regularization
    \item \textbf{Corruption prior} - Assume that some data points are corrupted or extreme, and can thus affect the final accuracy of our predictions - convert to percentiles, ranks
    \item \textbf{Low-Rank prior} - Assume that the actual variation only occurs in a small subset of the original predictors - use dimensionality reduction methods such as PCA
    \begin{itemize}
        \item In PCA, each principal component ($PC_i$) captures the $i$th direction with the most variance in the dataset.
    \end{itemize}
\end{itemize}
\end{document}
